\oldpage[IV-2]{187}
\begin{remarks}[7.1.12]
\label{IV.2.s7.1.12-fr}
\begin{enumerate}
\item[(i)] Se ignora si, recíprocamente, todo
anillo local noetheriano universalmente catenario
es siempre formalmente catenario.
\item[(ii)] Un anillo local noetheriano $A$
\emph{de dimensión} 1 es necesariamente
equidimensional, pues su ideal maximal no puede
ser minimal (ya que, en ese caso, $A$ tendría
dimensión 0). Como $\dim(\widehat A)=1$,
esto muestra que $A$ es incluso
\emph{formalmente equidimensional} y, \emph{a
fortiori}, formalmente catenario (y, por tanto,
\emph{universalmente catenario}). En cambio,
el anillo local de dimensión 2 definido en
(5.6.11), que es catenario pero no universalmente
catenario, \emph{a fortiori} no es formalmente
catenario.
\end{enumerate}
\end{remarks}
